% Generated draft: requires mathematical and source review.
\documentclass[11pt]{article}
\usepackage[margin=25mm]{geometry}
\usepackage{fontspec}
\setmainfont{Latin Modern Roman}
\usepackage{amsmath,amssymb,mathtools,unicode-math}
\setmathfont{Latin Modern Math}
\usepackage{enumitem,needspace,xurl,hyperref}
\setlength{\parindent}{0pt}
\setlength{\parskip}{5pt}
\setlength{\emergencystretch}{3em}
\newcommand{\N}{\mathbb{N}}
\newcommand{\Z}{\mathbb{Z}}
\newcommand{\Q}{\mathbb{Q}}
\newcommand{\R}{\mathbb{R}}
\newcommand{\C}{\mathbb{C}}
\newcommand{\F}{\mathbb{F}}
\newcommand{\A}{\mathbb{A}}
\newcommand{\PP}{\mathbb{P}}
\newcommand{\E}{\mathbb{E}}
\newcommand{\eps}{\varepsilon}
\newcommand{\dd}{\,\mathrm{d}}
\begin{document}
\section*{UnsolvedMath ID: 701022}
\textbf{Source problem code:} catalogue entry
\subsection{Definitions and mathematical statement}
Let \(n,m\geq 1\) be integers, let \(K\) be an infinite field of characteristic \(p>0\), and put \(M(n)=\operatorname{Mat}_{n\times n}(K)\). Let \(G(n)\leq\operatorname{GL}_n(K)\) be one of the classical matrix groups described in the source as ``special, orthogonal, or symplectic,'' subject to the unresolved specifications below. Its diagonal conjugation action on \(M(n)^m\) is
\[
g\cdot(A_1,\ldots,A_m)=(gA_1g^{-1},\ldots,gA_mg^{-1}).
\]
The coordinate algebra is
\[
R=K[M(n)^m]=K[x_{ij}^{(r)}\mid 1\leq i,j\leq n,\ 1\leq r\leq m],
\]
and its invariant subalgebra is
\[
R^{G(n)}=\{f\in R\mid f(gA_1g^{-1},\ldots,gA_mg^{-1})=f(A_1,\ldots,A_m)\text{ for every }g\in G(n)\text{ and }(A_1,\ldots,A_m)\in M(n)^m\}.
\]
\textbf{Problem.} For the intended choices of \(G(n)\), for every admissible \(n\), every \(m\geq1\), and every infinite field \(K\) of positive characteristic, describe explicit generators \(S\subseteq R^{G(n)}\) such that
\[
R^{G(n)}=K[S].
\]
\textbf{Unresolved specifications.} The supplied wording does not expand ``special'' into a precise group name; its subsequent comments explicitly discuss special orthogonal groups. It also does not specify the forms defining the orthogonal and symplectic groups or the convention for orthogonal groups in characteristic \(2\). These choices must be supplied to obtain a fully specified family of invariant algebras. Under the usual definition using a nondegenerate alternating form, a symplectic group requires \(n\) to be even.

\textbf{Supplied historical qualifications.} The source attributes the characteristic-zero generators to C. Procesi (1976). Its comment of 2001 reports a positive-characteristic solution except for orthogonal groups in characteristic \(2\) and special orthogonal groups of even degree, citing A. N. Zubkov (1999). Its comment of 2009 additionally reports a solution for special orthogonal groups of even degree over infinite fields of odd characteristic, citing A. A. Lopatin (2009). These are supplied historical assessments; they do not establish the current status of any remaining case.
\subsection{Short English statement}
Find explicit polynomial generators for the invariants of simultaneous conjugation on \(m\) square matrices over an infinite field of positive characteristic, for the classical groups intended by the source. The exact group families and defining forms require clarification. The supplied 2001 and 2009 comments qualify which cases were reported solved; the registry's \texttt{open} label is not an independently verified status.
\subsection{Sources}
\begin{itemize}
\item Local UnsolvedMath registry, permanent ID \(701022\), exact source code \texttt{catalogue entry}, title \emph{Kourovka Notebook Problem 14.37}. The supplied record labels the problem \texttt{open} and was updated on \texttt{2026-10-16T14:00:00Z}; this is registry metadata, not an independently verified literature assessment.
\item \emph{The Kourovka Notebook (Unsolved Problems in Group Theory)}, original problem 14.37, 14th issue (1999), proposed by A. N. Zubkov. The supplied background identifies the consulted edition as the 21st issue edition (2026), \texttt{arXiv:1401.0300v48}. The supplied collection description names E. I. Khukhro and V. D. Mazurov as editors. These source identifications have not been independently checked.
\item C. Procesi, \emph{Adv. Math.} \textbf{19} (1976), 306--381. The source cites this for generators in characteristic \(0\); no article title is supplied.
\item A. N. Zubkov, \emph{Algebra and Logic} \textbf{38}, no. 5 (1999), 299--318. The source's comment of 2001 cites this for the positive-characteristic solution with the stated exceptions: orthogonal groups in characteristic \(2\) and special orthogonal groups of even degree. No article title is supplied.
\item A. A. Lopatin, \emph{J. Algebra} \textbf{321} (2009), 1079--1106. The source's comment of 2009 cites this for special orthogonal groups of even degree over infinite fields of odd characteristic. No article title is supplied.
\end{itemize}
\end{document}
